\section{策略梯度直接优化动作分布}
\label{sec:16-Control-and-Reinforcement-Learning/02-Reinforcement-Learning/09}

\subsection{值方法在动作空间面前的两难}

价值方法先估计每个动作的长期回报，再在每个状态上选值最大的动作。这套逻辑在离散动作中很顺畅：4 个按钮，4 个 Q 值，取 \(\argmax\) 就是比大小。

但动作空间一旦变成连续的——机器人 7 自由度的关节角度、自动驾驶的方向盘转角——每一步都要在高维连续空间上做 \(\argmax\)。这是个嵌在决策循环里的非凸优化子问题，Q 函数本身还在不断随训练变化。值方法没有直接给你策略，它给你一堆数，然后让你在每个状态下临时解一个优化问题。

还有另一种情况：你本身就想要一个随机策略。博弈对抗中确定性策略可能被对手完美利用；探索密集的环境中，softmax 输出的概率分布比单点的 \(\argmax\) 包含更多信息。

策略梯度的思路很直接：与其绕道值函数再取 \(\argmax\)，不如直接参数化策略

\[
\pi_\theta(a \mid s)
\]

然后调整参数 \(\theta\)，让高回报的轨迹变得更可能，低回报的轨迹变得更不可能。把决策问题变成一个参数优化问题。

\subsection{第一步就撞到的麻烦}

假设轨迹

\[
\tau = (s_0, a_0, r_1, s_1, a_1, r_2, \ldots, s_T)
\]

在初始分布 \(\rho_0\) 和环境转移 \(P\) 下发生。轨迹的联合概率是

\[
p_\theta(\tau) = \rho_0(s_0) \prod_{t=0}^{T-1}
\pi_\theta(a_t \mid s_t) \, P(s_{t+1} \mid s_t, a_t).
\]

目标是期望回报

\[
J(\theta) = \E_{\tau \sim p_\theta}[G(\tau)].
\]

想对 \(\theta\) 求梯度，立刻遇到一个问题：\(\theta\) 不仅影响给定轨迹上回报的高低——它首先影响哪些轨迹会被采样到。\(\E_{\tau \sim p_\theta}\) 中的分布下标依赖 \(\theta\)。不能把一批采样轨迹当作与参数无关的固定数据，直接对回报取导——数据生成过程本身就含 \(\theta\)。

这和普通的监督学习参数优化不一样。监督学习中，数据和模型参数独立——你用固定数据集算 loss，对参数求导。强化学习中，采到的数据是什么，本身取决于策略参数当下的值。梯度必须穿过这个依赖。

\subsection{log-derivative：把分布求导变成采样平均}

恒等式

\[
\nabla_\theta p_\theta(\tau) = p_\theta(\tau) \, \nabla_\theta \log p_\theta(\tau)
\]

是最关键的一步。它把对概率密度求导变成了密度自身乘上一个对数导数。代入积分号：

\[
\begin{aligned}
\nabla_\theta J(\theta)
&= \nabla_\theta \int G(\tau) \, p_\theta(\tau) \, d\tau \\
&= \int G(\tau) \, \nabla_\theta p_\theta(\tau) \, d\tau \\
&= \int G(\tau) \, p_\theta(\tau) \, \nabla_\theta \log p_\theta(\tau) \, d\tau \\
&= \E_{\tau \sim p_\theta}
\Bigl[ G(\tau) \, \nabla_\theta \log p_\theta(\tau) \Bigr].
\end{aligned}
\]

现在它是一个期望了——可以采样。环境转移 \(P\) 和初始分布 \(\rho_0\) 都不含 \(\theta\)，log-概率的梯度只剩策略部分：

\[
\nabla_\theta \log p_\theta(\tau) = \sum_{t=0}^{T-1}
\nabla_\theta \log \pi_\theta(a_t \mid s_t).
\]

一条轨迹的一次采样就给出了无偏 Monte Carlo 估计。直觉清晰得不需要公式也能说：回报高的轨迹，把上面的每个动作概率往上拉；回报低的，往下压。

\subsection{一个小例子：两臂赌博机}

一条臂以概率 \(p_1\) 给奖 1，另一条以概率 \(p_2\) 给奖 1，每次只能选一条。策略参数化为 \(\pi_\theta(\text{左}) = \theta\)。采样 100 次，左侧选中 40 次，其中 8 次得奖；右侧选中 60 次，其中 30 次得奖。

平均回报约 \(0.38\)。REINFORCE 会把左侧期望回报高的动作概率往上拉。在这里，右侧回报 \(30/60 = 0.5\)，左侧 \(8/40 = 0.2\)，所以梯度会降低 \(\theta\)（减少选左侧的概率）。

如果没有 log-derivative 恒等式，你可能会想：回报对参数的导数怎么算？选左和选右的概率加起来为 1，调参数改变概率分配，但回报是随机变量。log-derivative 把``改变分布''的效果翻译成了``在样本上加权重''的期望——权重就是回报本身。

\subsection{去掉过去：reward-to-go}

完整 \(G(\tau)\) 乘在轨迹上每个时间步的 score 上，意味着动作 \(A_t\) 的梯度被它发生之前的奖励加权了。当前动作改变不了过去——这部分加权在期望意义下是零，但采样上引入额外波动。

利用条件期望的性质，可以只保留从当前时刻开始的未来回报：

\[
\nabla_\theta J(\theta)
= \E_\pi \Bigl[ \sum_t G_t \, \nabla_\theta \log \pi_\theta(A_t \mid S_t) \Bigr].
\]

期望不变，方差下降。REINFORCE 用样本 \(G_t\) 做更新：

\[
\theta \leftarrow \theta + \alpha \, G_t \, \nabla_\theta \log \pi_\theta(A_t \mid S_t).
\]

但注意：\(G_t\) 进入乘数，所以奖励的绝对尺度直接缩放有效步长。奖励放大十倍，相当于学习率放大十倍。奖励缩放、折扣约定和 episode 长度都绑在步长上，不是分离的超参数。

\subsection{baseline：减掉一个常数，期望不变}

对任意只依赖状态的函数 \(b(s)\)，

\[
\E_{A \sim \pi_\theta(\cdot \mid s)}
\bigl[ b(s) \, \nabla_\theta \log \pi_\theta(A \mid s) \bigr]
= b(s) \, \nabla_\theta \sum_a \pi_\theta(a \mid s) = 0.
\]

所以把 \(G_t\) 换成 \(G_t - b(S_t)\)，梯度的期望不变。关键在方差：选 \(b(S_t)\) 接近状态价值 \(V^\pi(S_t)\)，差值就衡量``这个动作比该状态的平均水平好多少''，正负分得更开，方差大幅下降。

一个具体的数字直觉：如果一个小区域内所有动作的 \(G_t\) 都是正的（比如 50 到 60 之间），REINFORCE 会同时抬高所有动作概率，有效信号被公共的大正数淹没。减去一个接近均值（比如 55）的 baseline，就会把 55 附近的动作变为接近零的信号，远高于 55 的变正，远低于 55 的变负。两种信号分开了，梯度更干净。

但 baseline 若依赖当前动作，上述消去一般不成立，需要额外修正。实践中常用 batch 平均回报做 baseline——简单，但与每个样本有 batch 内相关性，统计行为不是完全白噪声。

\subsection{策略参数化的形式就是可探索空间的形状}

离散动作常用 softmax：

\[
\pi_\theta(a \mid s) = \frac{\exp z_\theta(s, a)}{\sum_{a'} \exp z_\theta(s, a')}.
\]

连续动作常用 Gaussian 策略，网络输出均值 \(\mu_\theta(s)\) 和标准差 \(\sigma_\theta(s)\)：

\[
\pi_\theta(a \mid s) = \frac{1}{\sqrt{2\pi}\sigma_\theta(s)}
\exp\Bigl( -\frac{(a - \mu_\theta(s))^2}{2\sigma_\theta^2(s)} \Bigr).
\]

标准差的演化直接影响探索。如果 \(\sigma\) 过早塌缩趋近于零，策略趋近确定性，未尝试过的动作区域几乎没有梯度信号，探索停止，可能困在局部最优。如果 \(\sigma\) 一直很大，采样回报的方差高，梯度噪声大。Entropy bonus 是延缓过早集中的一种常见手段，但它在目标函数中加入了一个新的偏好项——不再是在原问题上做纯粹优化，而是在``原目标 + 熵正则''的修正问题上做优化。目标和原始问题的差距值得留意。

\subsection{策略梯度定理给的是局部方向，不是路线图}

策略梯度定理确保梯度方向（期望意义下）指向期望回报增加的方向。在足够小的步长下，一直沿梯度方向走，期望回报不会下降。但回报曲面可能高度非凸——策略参数化叠加上环境动力学的非线性，曲面上满布鞍点和局部极值。策略梯度不保证走出它们。

一个极端但常见的场景：稀疏奖励的长时域任务。如果任务成功才给 \(+1\)，其余全为零，只有极少数轨迹碰巧触发奖励。绝大多数采样轨迹回报为零——梯度也为零。偶尔一条幸运轨迹击中奖励，给一个巨大的梯度脉冲。方差大到几乎不可用的程度。这不是学习率的问题——信号本身就有结构性缺失。

策略梯度的无偏性是它在理论上的核心优势——不会因为价值函数近似而引偏。但当方差大到连无偏估计的符号都不稳定时，无偏性在实用上的价值就打折扣了。Actor-Critic 的出发点就是用学习到的价值函数替代完整的 Monte Carlo 回报——接受偏差，换取方差。
